% Generated TeX body; no PDF or document-style claim.
% Maintained source: manuscripts/graph-backends/static-confluent-rounding-and-tail-bounds.md
% Recovered-source SHA-256: 1c1173216cdf8559e67776ba21896ddca9a02cb06d06610711d0632869ed13fd
% Rebuild: python3 tools/build_manuscripts.py
% Common packages: amsmath, amssymb, hyperref, longtable, booktabs, array.
% Shared original Subset Sum preamble: manuscripts/paper-style.tex.

\hypertarget{af-385fd70686-deterministic-confluent-rounding-implementation-moments-and-load-tails}{%
\section{Deterministic confluent rounding: implementation, moments, and load tails}\label{af-385fd70686-deterministic-confluent-rounding-implementation-moments-and-load-tails}}

\textbf{Ready for review.} Complete near-linear simulation and Parts B--D simultaneous completed-sink moment and tail bounds, including the budget proportional to actual demand. Part E proves the scoped local warm-start lemma.

\textbf{Dependencies and scope.} The precise published base-2 Conflt operations and deterministic dynamic-tree interface are identified. Reopening completed sinks needs separate dynamic replenishment accounting and is not used by the completed main matching theorem.

The \href{https://github.com/icekylinx/algorithm-frontier/blob/main/manuscripts/graph-backends/verify/README.md}{verification note} identifies executable coverage. This status does not claim external peer review or proof-assistant verification.

\textbf{Results.} Starting from a supplied splittable flow with positive nonsink demands, deterministic confluent rounding chooses a single outgoing route at every nonsink. Part A implements the rounding in \(O((n+m) polylog(n+m+2))\) exact arithmetic operations. Parts B--D give complete bounds that include sinks completed before the last step. In particular, if \(C\) bounds every initial node congestion and \(W\) is total actual demand, the final loads satisfy \[
C\sum_j\frac{2^{W_j/C}-1}{(1+W_j/C)^3}=O(W).
\] The left side dominates the total component weight up to an absolute constant. This stronger mass-sensitive bound does not require the number of sinks times \(C\) to be near-linear in the represented weight.

The proof package was already established in the 24 October 2026 checkpoint. Its near-linear simulation, global moment, load-tail potential, mass-sensitive potential, and local rounding lemma are assembled here without dropping their arguments. A dynamic scheme that repeatedly reopens completed trees needs its own accounting; that separate research question is stated precisely in Parts C--E. It is not a premise of the current \href{https://github.com/icekylinx/algorithm-frontier/blob/main/problems/general-matching/proof.md}{four-fifths maximum-matching theorem}, which does not use this package.

\hypertarget{af-385fd70686-theorem-map-and-model}{%
\subsection{Theorem map and model}\label{af-385fd70686-theorem-map-and-model}}

Let \(n,m\) count vertices and physical arc records of the supplied graph, including parallel records. Zero arcs are removed in the input scan. Nonsink demands are strictly positive; sink demands may be zero. Node congestion means own demand plus incoming ordinary flow, equivalently outgoing flow at a nonsink. Original sinks have no ordinary outgoing arcs. A confluent output is a forest of original arc witnesses, one outgoing arc per original nonsink, reaching original sinks.

\begin{longtable}[]{@{}
  >{\raggedright\arraybackslash}p{(\columnwidth - 4\tabcolsep) * \real{0.3333}}
  >{\raggedright\arraybackslash}p{(\columnwidth - 4\tabcolsep) * \real{0.3333}}
  >{\raggedright\arraybackslash}p{(\columnwidth - 4\tabcolsep) * \real{0.3333}}@{}}
\toprule\noalign{}
\begin{minipage}[b]{\linewidth}\raggedright
Part
\end{minipage} & \begin{minipage}[b]{\linewidth}\raggedright
Proved conclusion
\end{minipage} & \begin{minipage}[b]{\linewidth}\raggedright
Initial assumptions
\end{minipage} \\
\midrule\noalign{}
\endhead
\bottomrule\noalign{}
\endlastfoot
\protect\hyperlink{af-385fd70686-implementation}{A} & Deterministic near-linear rounding; maximum sink load at most \(C(1+log_2 F)\) & Supplied flow with all initial node loads at most \(C\); \(F≥1\) sinks \\
\protect\hyperlink{af-385fd70686-exponential-moment}{B} & Sum over \textbf{all} final sinks of \(2^x\) is at most \(P(3+½log_2 P)\) & Nonsink congestion at most \(C\); initial active potential \(P=Σ2^x\) \\
\protect\hyperlink{af-385fd70686-load-tail}{C} & Number of final loads at least \(τC\) is at most \(F(τ+3)/2^(τ+1)\) for \(τ≥1\) & All initial node loads at most \(C\) \\
\protect\hyperlink{af-385fd70686-mass-sensitive-tail}{D} & Number is at most \((τ+5)(W/C)/(7·2^(τ-1))\) for \(τ≥4\); total weighted budget \(O(W)\) & Exact total demand \(W\), after trimming if necessary \\
\protect\hyperlink{af-385fd70686-local-warm-start}{E} & Round a chosen region while preserving exterior assignments, with load controlled by its actual initial potential & Dirty rows supplied; exterior trees represented as sink anchors with their actual demands \\
\end{longtable}

Empty graphs are immediate. A nonempty valid flow with positive nonsink demand must have a sink; in Parts A--C the nonempty initial sink set gives \(P≥1\). If \(W=0\), Part D is the trivial all-zero statement. Section and equation numbers are local to each lettered part.

\hypertarget{af-385fd70686-named-dependencies-and-what-is-proved-here}{%
\subsection{Named dependencies and what is proved here}\label{af-385fd70686-named-dependencies-and-what-is-proved-here}}

The rounding operations and the base congestion theorem are those of Chen, Kleinberg, Lovász, Rajaraman, Sundaram, and Vetta, \href{https://www.cs.cornell.edu/~rdk/papers/conflu.pdf}{\emph{(Almost) Tight Bounds and Existence Theorems for Single-Commodity Confluent Flows}}, \textbf{Section 4.1, Theorem 3, Lemmas 4--6}. This page uses that section's base-2 algorithm, not the different refinement in Section 4.2. The original source gives a polynomial-time implementation; Part A proves the near-linear simulation of its permitted operations.

For data structures, use the deterministic dynamic-forest interface of Alstrup, Holm, de Lichtenberg, and Thorup, \href{https://arxiv.org/abs/cs/0310065}{\emph{Maintaining Information in Fully-Dynamic Trees with Top Trees}}. The required summaries and update charges are specified in Part A. The trimming operation has a complete \href{https://github.com/icekylinx/algorithm-frontier/blob/main/manuscripts/graph-backends/exact-demand-flow-trimming.md}{local proof}, so no maximum-flow oracle is needed for that preprocessing.

Arithmetic costs are separate from operation counts. For integral data with every initial node congestion at most \(C\), ordinary arc values remain at most \(C\) and stored demand sums are at most \(nC\); \(O(log(nC+2))\)-bit values suffice. For the local version with large exterior demands, charge their actual input and sum bit lengths. The algorithm uses only integer/rational additions, subtractions, minima, and comparisons. Exponentials and logarithms occur only in the proof.

\hypertarget{af-385fd70686-part-a.-deterministic-implementation-of-the-rounding-operations}{%
\subsection{Part A. Deterministic implementation of the rounding operations}\label{af-385fd70686-part-a.-deterministic-implementation-of-the-rounding-operations}}

\hypertarget{af-385fd70686-a.1.-flow-model-and-permitted-operations}{%
\subsubsection{A.1. Flow model and permitted operations}\label{af-385fd70686-a.1.-flow-model-and-permitted-operations}}

Every nonsink has a strictly positive demand. The current graph has positive ordinary arc flows, and active sinks have no ordinary outgoing arcs. For each ordinary arc \texttt{u\ -\textgreater{}\ s} into an active sink, permit the reverse traversal \texttt{s\ -\textgreater{}\ u} as well. Regard the physical edge underlying this pair as bidirectional; all other ordinary edges remain one-way.

The Conflt operations are:

\begin{enumerate}
\def\labelenumi{\arabic{enumi}.}
\tightlist
\item
  If a nonsink has all its ordinary outgoing edges directed into one sink, aggregate it into that sink, mark one genuine outgoing edge, and transfer its demand to the sink.
\item
  On a directed cycle in the augmented graph, decrease flows along ordinary forward traversals and increase flows along reverse sink-frontier traversals by the minimum flow on a forward traversal. Delete zero-flow ordinary edges and their reverse traversals. The selected cycle must not consist of going along and immediately back over the same physical edge.
\item
  Let sink \(s_j\) have exactly one frontier neighbor \(v\), and let \(s_l\) be another sink neighbor of \(v\). Write \(a=f(v,s_j)\), \(b=f(v,s_l)\), and let \(s,t\) be the two sink loads. If \(s+2b<t\), redirect the entire \(b\) units from \(s_l\) to \(s_j\) and delete the former edge. Otherwise redirect the entire \(a\) units from \(s_j\) to \(s_l\), delete the former edge, and deactivate \(s_j\). Its completed load is \(s-a\), and the recipient load becomes \(t+a\). The comparison is exact and uses no exponential evaluation.
\end{enumerate}

For completeness, the potential invariant follows directly from these rules. Normalize flows and loads by the bound \(C\) on nonsink congestion. Aggregation preserves sink load, and cycle cancellation preserves sink load and never increases nonsink congestion. In the nondeactivating transfer, the inequality \(s+b<t-b\) makes the transfer from the larger load to the smaller load decrease the sum of any convex function of the two loads. In a deactivation, \(0<a≤1\), \(a+b≤1\), and \(t≤s+2b\). With \(q=2^a\) and \(r=2^(t-s)\), \[
r(q-1)\le \frac{4(q-1)}{q^2}\le1\qquad(1\le q\le2).
\] Consequently \(2^(t+a)-2^t≤2^s\), so the active exponential sum does not increase. A deactivated sink's final exponential is at most its pre-deactivation exponential, and hence at most the preceding active sum. Every marked route ends at an original sink, as proved in Section A.4. Thus initial sink loads at most \(C\) and \(F≥1\) sinks imply final congestion at most \(C(1+log_2 F)\). These are the source paper's Section 4.1 operations and guarantee, with their analysis included here; no stronger approximation guarantee is assumed. Ordinary directed cycles can also be canceled: the same flow update preserves conservation and decreases node congestion. Thus a separate initial acyclic-flow conversion can be used, but the local simulation below does not depend on choosing a globally cycle-free support at each iteration.

Self-loops can be canceled immediately. Parallel ordinary arcs may either be coalesced, retaining one original physical edge as the eventual marking witness, or be kept with distinct physical edge identifiers. In the latter implementation, a two-edge cycle using \textbf{different} physical edges is a valid cancellation: it redirects flow between those edges and permanently deletes at least one. Only the forward/reverse pair of the \textbf{same} physical edge is excluded. There are only linearly many initial arcs and retarget/coalescence events.

\hypertarget{af-385fd70686-a.2.-a-mixed-selected-forest}{%
\subsubsection{A.2. A mixed selected forest}\label{af-385fd70686-a.2.-a-mixed-selected-forest}}

Maintain an undirected forest \texttt{F} consisting of a subset of the current physical edges. Each one-way selected edge points toward the root of its selected tree. A selected bidirectional edge may be traversed in either direction.

Equivalently, contract each connected component of the selected bidirectional edges. The selected one-way edges form a directed forest with outdegree at most one. Every full selected tree therefore has a unique terminal \textbf{root component}. Its vertices are connected using selected bidirectional edges only; all other selected edges point toward it.

Initially \texttt{F} is empty, so the invariant holds. Keep a designated full-tree root inside its root component.

Maintain a second dynamic forest containing only the selected bidirectional edges. At each vertex store its allowed outgoing arcs whose underlying physical edge is \textbf{not} selected in \texttt{F}. Aggregate the existence of such a candidate over each bidirectional component. Both directions of an already selected physical edge are excluded; this avoids a vacuous two-step cycle.

\hypertarget{af-385fd70686-a.3.-a-local-simulation-step}{%
\subsubsection{A.3. A local simulation step}\label{af-385fd70686-a.3.-a-local-simulation-step}}

Choose a live selected tree and its root component \texttt{Q}.

\hypertarget{af-385fd70686-case-a-an-unselected-outgoing-arc-exists-at-a-vertex-of-q}{%
\paragraph{Case A: an unselected outgoing arc exists at a vertex of Q}\label{af-385fd70686-case-a-an-unselected-outgoing-arc-exists-at-a-vertex-of-q}}

Take any allowed such traversal \texttt{u\ -\textgreater{}\ v}, and reroot the selected tree at \texttt{u}. This reroot uses only bidirectional selected edges, so no one-way edge is reversed.

If \texttt{u,v} are in different selected trees, insert their physical edge into \texttt{F}, directed from the former root component toward \texttt{v}. This preserves the mixed-forest invariant.

If they are already in one selected tree, the selected path from \texttt{v} to \texttt{u} respects all one-way orientations: it follows the directed component forest toward \texttt{Q}, then uses bidirectional edges within \texttt{Q}. The new traversal followed by that path is therefore a genuine directed augmented cycle. Cancel it by operation 2. At least one forward ordinary edge reaches zero. Delete every zero edge. If the new edge remains positive, put it into the selected forest after the zero edges have broken the cycle.

\hypertarget{af-385fd70686-case-b-no-unselected-outgoing-arc-exists-at-q}{%
\paragraph{Case B: no unselected outgoing arc exists at Q}\label{af-385fd70686-case-b-no-unselected-outgoing-arc-exists-at-q}}

If \texttt{Q} is an isolated active sink with no incident edge, finish that sink's component.

A singleton nonsink with no outgoing edge is impossible while its positive demand is conserved. Otherwise \texttt{Q} contains sinks and their frontier neighbors, with selected bidirectional edges forming a tree. Every incoming ordinary edge of a sink in \texttt{Q} must already be selected in this tree: otherwise its reverse would be an eligible unselected outgoing traversal. Every outgoing ordinary edge of a frontier in \texttt{Q} is likewise selected inside \texttt{Q}: otherwise it would be eligible. Therefore \texttt{Q} is the entire local sink/frontier incidence tree, although nonsink frontier nodes may have incoming arcs from outside \texttt{Q}.

If this tree has a leaf frontier, that frontier has only one ordinary outgoing neighbor, so perform operation 1. If it has no leaf frontier, it has a leaf sink. That sink has one frontier neighbor, and that frontier has another sink neighbor; perform operation 3. These are exactly valid source-paper operations.

Thus every processed root component yields a selected link, a positive-flow edge deletion, an aggregation, a sink deactivation, or completion of an isolated sink. There is no unsuccessful global search.

\hypertarget{af-385fd70686-a.4.-dynamic-tree-primitives}{%
\subsubsection{A.4. Dynamic-tree primitives}\label{af-385fd70686-a.4.-dynamic-tree-primitives}}

Use a full dynamic forest for connectivity, roots, link/cut, and path flow operations, and an Euler-tour or top-tree representation of the selected bidirectional subforest for component aggregates and candidate retrieval.

For a directed augmented cycle, the ordinary flow value decreases on forward traversals and increases on reverse traversals. A link-cut path aggregate stores the minimum ordinary flow among forward-oriented edges, and separate forward/reverse summaries. Lazy signed path additions and path reversals update these summaries. Reverse traversals have no saturation upper bound. The bottleneck is the minimum among the forward ordinary traversals, including the newly proposed nonforest traversal when it is forward. After the update, expose the affected path and enumerate all zero-flow forward edge identifiers before making any cut. Temporarily exclude each already reported edge from the minimum summary, then cut the recorded edges. Each extraction and deletion is charged to an original edge's permanent removal; no later search assumes that the cut path is still connected.

An edge that becomes incident into a sink acquires its reverse traversal once. If that edge is selected, insert it into the bidirectional subforest. The former two bidirectional components were joined by a directed selected edge, so their union still has at most one outgoing selected one-way edge.

On aggregation of a frontier \texttt{v} into a sink \texttt{s}, contract the selected edge \texttt{v-s}, retarget incoming arcs of \texttt{v} to \texttt{s}, and delete its outgoing arcs to \texttt{s}. Incoming selected edges can be cut and relinked individually. Every original arc is retargeted at its head at most once: its head changes from an ordinary vertex to a sink, and active sinks are never aggregated into other sinks. A sink deactivated later has no remaining incoming arc. Thus these adjacency scans sum to \texttt{O(m)}.

The source-paper sink-deactivation comparison reads two sink congestions and two arc flows, performs one flow transfer, and removes an edge. All needed quantities support point access. Cycle cancellations leave the congestion of every active sink unchanged, so they do not require enumerating all touched sink vertices. No exponential is evaluated by the algorithm: the exponential potential is only an analysis device, while the deactivation decision uses ordinary exact additions and a comparison. Integer supplied flows and demands remain integral throughout.

Every marked aggregation edge retains its \textbf{original head vertex}, rather than only the current contracted sink representative. That original head was either an original sink or a vertex aggregated earlier. Hence the marked edges form a genuine forest in the original graph, and routing them reaches the correct original sink. Selected bidirectional components containing actual sink vertices may occur in nonroot branches; the algorithmic root component means the unique terminal component after contracting selected bidirectional edges, not the set of all actual sinks in a selected tree.

\hypertarget{af-385fd70686-a.5.-amortized-count}{%
\subsubsection{A.5. Amortized count}\label{af-385fd70686-a.5.-amortized-count}}

There are \texttt{O(n+m)} permanent ordinary-edge removals, vertex aggregations and sink completions. A cycle cancellation permanently removes at least one positive-flow edge, so there are \texttt{O(m)} such cancellations.

A selected link between different trees reduces the number of selected forest components by one. Apart from these links, the number of components can increase only when selected edges are cut because of a permanent edge deletion or the linearly many explicit retarget/coalescence operations. Hence the number of pure selected links is also \texttt{O(n+m)}. Insertion of a surviving cycle edge is charged to its cycle cancellation.

For leaf selection in Case B, maintain a component minimum over the selected bidirectional degrees and the sink/frontier type: a tree with at least one edge has a marked leaf of one of the required types. Degree and candidate marks change only at endpoints of the linearly many membership or retarget operations. Completed isolated sinks need no further scheduling, but are kept in the analysis potential until termination unless the explicit deactivation rule removed them. This distinction prevents an uncharged disappearance of completed mass.

Each forest membership change affects only the two endpoint candidate lists. Edge retargeting, reverse activation and degree changes also occur linearly many times. With logarithmic dynamic-tree overhead, the total bound is \texttt{O((n+m)\ polylog\ n)} arithmetic operations, plus the bit complexity of the supplied flow values.

\hypertarget{af-385fd70686-integer-word-bounds}{%
\paragraph{Integer word bounds}\label{af-385fd70686-integer-word-bounds}}

For the matching-state application, the initial ordinary flows and demands are nonnegative integers, the nonsink demands are positive, and every initial node congestion is at most the integer \texttt{L}. All subsequent changes use additions, subtractions, minimum operations and the exact linear sink comparison. Consequently integrality is preserved. An ordinary arc's flow is at most the total outgoing flow of its nonsink tail, whose congestion never increases and is therefore at most \texttt{L}. Sink demand is a sum of original demands; even the crude bound \texttt{nL} suffices for intermediate stored quantities. Thus \texttt{O(log(nL))} bits per stored arithmetic value are enough. When \texttt{L} is polynomial in the input size, these fit the logarithmic-word RAM used by the matching-state framework. No irrational exponential or logarithm is evaluated during execution.

\hypertarget{af-385fd70686-a.6.-application-to-a-bottom-allocation-ledger}{%
\subsubsection{A.6. Application to a bottom allocation ledger}\label{af-385fd70686-a.6.-application-to-a-bottom-allocation-ledger}}

Fix an injection \texttt{mu}. For each assigned bin \texttt{v}, its assigned row has total allocation \texttt{L}; send that allocation from \texttt{v} to the row's bin neighbors. Remove self-loops. Before self-loop removal, the supply at an assigned bin is \texttt{L-inflow(v)}, at least its weight. A free bin is a sink; give it its own weight as local demand. Its total demand plus incoming flow is at most \texttt{L}. The supplied flow therefore has node congestion at most \texttt{L} and routes at least the original bin weights.

The confluent rounding selects one genuine outgoing row-neighbor edge for the row assigned to each assigned bin, reaches the original free bins, and forms a row-labelled witness forest. The original assignment of each row to its tail bin is retained; only its other selected neighbor is chosen by rounding. Its tree weights are at most \texttt{L(1+log\_2\ \textbar{}F\_free\textbar{})}. This is a deterministic static balanced-forest construction in near-linear time under the dynamic-tree interface proved above.

This does not solve the dynamic matching task. After a gain, a locally repaired affected tree can attach mass to surviving trees and increase their weight. A charged scheme maintaining the balance bound throughout expansions, contractions, gains and peels remains necessary.

\hypertarget{af-385fd70686-part-b.-an-exponential-moment-including-every-completed-sink}{%
\subsection{Part B. An exponential moment including every completed sink}\label{af-385fd70686-part-b.-an-exponential-moment-including-every-completed-sink}}

\hypertarget{af-385fd70686-b.1.-statement}{%
\subsubsection{B.1. Statement}\label{af-385fd70686-b.1.-statement}}

Suppose Conflt starts from a splittable flow with nonsink congestion at most \texttt{C\textgreater{}0} and nonnegative sink loads. Let

\[
P=\sum_{\mathrm{initial\ sinks}\ s}2^{b_s/C}.
\]

Then every final confluent tree load \texttt{B\_s} satisfies the published maximum bound \texttt{B\_s\textless{}=C\ log\_2\ P}. In addition,

\[
\sum_{\mathrm{final\ sinks}\ s}2^{B_s/C}\le P+\frac P2(4+\log_2P)=P\left(3+\frac{\log_2P}{2}\right).\tag{1}
\]

It is harmless to include original sinks whose final tree contains only their own demand. In the usual case \texttt{b\_s\textless{}=C} initially and there are \texttt{F} sinks, \texttt{P\textless{}=2F}, so (1) is \texttt{O(F\ log(F+2))}.

\hypertarget{af-385fd70686-b.2.-a-potential-that-also-pays-for-completed-sinks}{%
\subsubsection{B.2. A potential that also pays for completed sinks}\label{af-385fd70686-b.2.-a-potential-that-also-pays-for-completed-sinks}}

Normalize loads by \texttt{C}, writing \texttt{x\_s=b\_s/C}. The published active-sink potential is

\[
\Phi=\sum_{\mathrm{active\ sinks}}2^x\le P.
\]

Put \texttt{D=4+log\_2\ P} and introduce

\[
\Psi=\sum_{\mathrm{active\ sinks}}H(x),\qquad H(x)=(D-x)2^x.
\]

All active loads lie in \texttt{{[}0,log\_2\ P{]}}. On that interval,

\[
\begin{aligned}H'(x)&=2^x[(D-x)\ln2-1]>0,\\H''(x)&=2^x[(D-x)(\ln2)^2-2\ln2]>0.\end{aligned}
\]

Thus \texttt{H} is increasing and convex throughout the range used by Conflt. Cycle cancellations and node aggregations preserve active sink loads. In the nondeactivating balancing case, flow moves from the higher sink load to the lower without crossing; convexity makes \texttt{Psi} nonincreasing.

It remains to charge a deactivation. Denote the two current normalized sink loads by \texttt{s,t}, and the normalized frontier flows to them by \texttt{a,b}. Then \texttt{a\textgreater{}0}, \texttt{a+b\textless{}=1}, and the deactivation case satisfies

\[
s+b\ge t-b.
\]

The frontier's \texttt{a} units are moved from sink \texttt{s} to sink \texttt{t}; the former sink is deactivated with final load \texttt{s-a}, and the surviving sink load becomes \texttt{t+a}.

Define

\[
A=2^s,\qquad q=2^a,\qquad r=2^{t-s},\qquad u=r(q-1).
\]

Here \texttt{1\textless{}q\textless{}=2}. The source's active-potential inequality gives \texttt{u\textless{}=1}. For completeness, it follows directly from

\[
r\le2^{2b}\le2^{2(1-a)},\qquad r(q-1)\le4(2^{-a}-2^{-2a})\le1.
\]

The decrease in \texttt{Psi}, divided by \texttt{A}, is exactly

\[
(D-s)(1-u)+u\log_2r+arq.
\]

Since \texttt{D-s\textgreater{}=4}, this is at least

\[
4+u[\log_2u+z(a)-4],\qquad z(a)=-\log_2(q-1)+\frac{aq}{q-1}.
\]

As a function of \texttt{q},

\[
z(q)=\frac{q\log_2q}{q-1}-\log_2(q-1),\qquad z'(q)=-\frac{\log_2q}{(q-1)^2}<0.
\]

Therefore \texttt{z(q)\textgreater{}=z(2)=2}. We obtain

\[
\frac{\Psi_{\mathrm{before}}-\Psi_{\mathrm{after}}}{A}\ge4+u\log_2u-2u\ge2,\qquad0\le u\le1.\tag{2}
\]

The final inequality holds because the displayed function is decreasing on \texttt{(0,1{]}}: its derivative is \texttt{log\_2\ u+1/ln\ 2-2\textless{}0}, and its value at one is two.

The deactivated sink's final exponential contribution is

\[
2^{s-a}=\frac Aq\le A.
\]

Thus (2) pays at least twice that contribution. Completion of an isolated active sink can either remain in the final active sum or be removed and charged directly, since its \texttt{H(x)\textgreater{}=4*2\^{}x}.

\hypertarget{af-385fd70686-b.3.-summing-the-charges}{%
\subsubsection{B.3. Summing the charges}\label{af-385fd70686-b.3.-summing-the-charges}}

The initial \texttt{Psi} is at most \texttt{(4+log\_2\ P)P}, since normalized initial loads are nonnegative. The potential is nonnegative. Consequently all deactivated final sinks together contribute at most \texttt{Psi\_initial/2}. The sinks remaining active at termination contribute at most \texttt{P} by the original potential. This proves (1).

If final component weights are only lower bounds on the routed demands, as in the matching-state bottom-ledger construction, replacing \texttt{B\_s} by those weights preserves (1).

\hypertarget{af-385fd70686-b.4.-consequence-for-balanced-row-forests}{%
\subsubsection{B.4. Consequence for balanced row forests}\label{af-385fd70686-b.4.-consequence-for-balanced-row-forests}}

In the bottom-ledger application \texttt{C=L}, initial sink loads are at most \texttt{L}, and a component's actual bin weight \texttt{W\_T} is at most its rounded guide-flow demand. Hence

\[
\sum_T2^{W_T/L}=O(F\log(F+2)).\tag{3}
\]

The physical carrier of a selected row tree has size \texttt{2W\_T-1}. Because \texttt{W\textless{}=O(L(2\^{}(W/L)-1))}, the budget

\[
L\sum_T(2^{W_T/L}-1)=O(LF\log(F+2)).
\]

dominates total carrier mass up to an absolute constant. This is not a claim about \texttt{sum\_T\ W\_T\ 2\^{}(W\_T/L)}. If a phase additionally guarantees \(F=O(K)\) and \(LK=O_a(N)\), this budget is \(O_a(N log N)\). A general optimum-gap phase need not have \(F=O(K)\): many bins may remain unmatched in every maximum matching. Part D removes that extra assumption by using total actual demand.

This is a genuine initial budget for an overflow/recruitment argument. It does not yet prove a dynamic budget: local rounding can rebuild potential at newly enlarged cores. A faster multilevel epoch theorem still needs to charge that replenishment explicitly.

\hypertarget{af-385fd70686-b.5.-source-and-attribution}{%
\subsubsection{B.5. Source and attribution}\label{af-385fd70686-b.5.-source-and-attribution}}

The Conflt operations and original active-sink potential are from Chen, Kleinberg, Lovasz, Rajaraman, Sundaram and Vetta, \emph{(Almost) Tight Bounds and Existence Theorems for Single-Commodity Confluent Flows}, Section 4.1, especially Lemma 6: \url{https://www.cs.cornell.edu/~rdk/papers/conflu.pdf}.

The auxiliary potential \texttt{Psi} and the simultaneous bound (1) are deductions already developed in the 24 October 2026 checkpoint. This restoration makes no new literature-priority claim.

\hypertarget{af-385fd70686-part-c.-a-tail-bound-in-terms-of-the-initial-number-of-sinks}{%
\subsection{Part C. A tail bound in terms of the initial number of sinks}\label{af-385fd70686-part-c.-a-tail-bound-in-terms-of-the-initial-number-of-sinks}}

\hypertarget{af-385fd70686-c.1.-the-lemma}{%
\subsubsection{C.1. The lemma}\label{af-385fd70686-c.1.-the-lemma}}

Suppose Conflt starts from a flow whose nonsink and sink congestions are at most \texttt{C}, with \texttt{F} sinks. Let \texttt{W\_j} denote the final demand at sink \texttt{j}, including the sinks deactivated before termination. For every real threshold \texttt{tau\ \textgreater{}=\ 1},

\[
\#\{j:W_j\ge \tau C\}
\le {F(\tau+3)\over 2^{\tau+1}}.                 \tag{1}
\]

In particular,

\[
\sum_{j:W_j\ge\tau C}W_j
=O\!\left({CF(\tau+1)^2\over2^\tau}\right).       \tag{2}
\]

The same bounds hold for actual selected-tree weights when the guide-flow demands dominate the bin weights. If the flow has first been trimmed to exact bin demand, the two quantities coincide.

\hypertarget{af-385fd70686-c.2.-a-capped-slope-potential}{%
\subsubsection{C.2. A capped-slope potential}\label{af-385fd70686-c.2.-a-capped-slope-potential}}

Normalize loads by \texttt{C}. Set \texttt{E=2\^{}tau} and

\[
\gamma=4\ln 2-1\in(0,4).
\]

Use the continuously differentiable function

\[
H_\tau(x)=
\begin{cases}
 (\tau+4-x)2^x,&0\le x\le\tau,\\
 E\,[4+\gamma(x-\tau)],&x\ge\tau.
\end{cases}                                      \tag{3}
\]

It is nonnegative, increasing, and convex. Below \texttt{tau}, its second derivative has the sign of

\[
(\tau+4-x)(\ln2)^2-2\ln2.
\]

which is positive because \texttt{tau+4-x\ \textgreater{}=\ 4}. Above \texttt{tau}, the function is linear, with the matching derivative \texttt{gamma\ E}. Its derivative is everywhere at most \texttt{gamma\ E}.

Let \texttt{Psi} be the sum of (3) over the currently active sinks. Node aggregation and cycle cancellation preserve sink congestions. The nondeactivating transfer in Conflt moves flow from a larger load to a smaller load without reversing their order; convexity therefore makes \texttt{Psi} nonincreasing.

\hypertarget{af-385fd70686-c.3.-the-deactivation-calculation}{%
\subsubsection{C.3. The deactivation calculation}\label{af-385fd70686-c.3.-the-deactivation-calculation}}

Write \texttt{s,t} for the two preoperation normalized sink loads, and \texttt{a,b} for the flows from the common frontier into these sinks. The deactivation branch has

\[
0<a\le1,\qquad a+b\le1,\qquad t\le s+2b.
\]

It completes the first sink at load \texttt{s-a}, and changes the surviving sink from \texttt{t} to \texttt{t+a}. The active-potential drop is

\[
\Delta=H_\tau(s)-[H_\tau(t+a)-H_\tau(t)].          \tag{4}
\]

This drop is always nonnegative.

If \texttt{s\ \textgreater{}=\ tau}, the first term is at least \texttt{4E} while the bracket is at most \texttt{a\ gamma\ E\ \textless{}=\ gamma\ E\ \textless{}\ 4E}.

Now suppose \texttt{s\ \textless{}\ tau}. If \texttt{t+a\ \textless{}=\ tau}, every value in (4) uses the exponential part of (3). The calculation in \protect\hyperlink{af-385fd70686-exponential-moment}{Part B}, with \texttt{D=tau+4}, applies because \texttt{D-s\ \textgreater{}=\ 4}. It gives \texttt{Delta\ \textgreater{}=\ 2*2\^{}s\ \textgreater{}\ 0}.

Finally, suppose \texttt{s\ \textless{}\ tau\ \textless{}\ t+a}. Since

\[
t+a\le s+2-a.
\]

we have \texttt{s\ \textgreater{}\ tau-2+a}. If the comparison point \texttt{tau-2+a} is negative, extend the exponential formula to that point solely for this inequality; it is still increasing there. Monotonicity and the derivative bound give

\[
{\Delta\over E}
> (6-a)2^{a-2}-a\gamma.                           \tag{5}
\]

The right side is decreasing on \texttt{0\ \textless{}=\ a\ \textless{}=\ 1}. Indeed the derivative of its first term is increasing on that interval, and even its value at \texttt{a=1}, after subtracting \texttt{gamma}, is

\[
\frac{5\ln2-1}{2}-(4\ln2-1)=\frac12-\frac32\ln2<0.
\]

Thus (5) is at least its value at one,

\[
\frac52-\gamma=\frac72-4\ln2>0.
\]

This proves the nonincrease for every deactivation.

There is a stronger charge when the completed load is at least \texttt{tau}. In that case \texttt{s-a\ \textgreater{}=\ tau}, and linearity on the tail gives

\[
\Delta\ge H_\tau(s)-a\gamma E
         =H_\tau(s-a)\ge4E.                       \tag{6}
\]

\hypertarget{af-385fd70686-c.4.-charge-all-completed-heavy-sinks}{%
\subsubsection{C.4. Charge all completed heavy sinks}\label{af-385fd70686-c.4.-charge-all-completed-heavy-sinks}}

Initially every normalized load is at most one. Because \texttt{tau\ \textgreater{}=\ 1},

\[
\Psi_{\rm initial}\le2F(\tau+3).                 \tag{7}
\]

Charge \texttt{H\_tau(final\ load)} for a deactivated sink only if that final load is at least \texttt{tau}. Equation (6) pays each such charge using an irreversible active-potential drop. At termination, every remaining heavy active sink is already present in the residual potential. Therefore

\[
\sum_{j:W_j/C\ge\tau}H_\tau(W_j/C)
\le 2F(\tau+3).                                  \tag{8}
\]

Equation (1) follows from \texttt{H\_tau(x)\textgreater{}=4E} on the tail. Since for \texttt{x\textgreater{}=tau}

\[
x\le \max\{\tau/4,1/\gamma\}
            [4+\gamma(x-\tau)],
\]

equation (2) follows as well.

\hypertarget{af-385fd70686-a-near-linear-weighted-exponential-budget}{%
\paragraph{A near-linear weighted exponential budget}\label{af-385fd70686-a-near-linear-weighted-exponential-budget}}

A useful consequence is

\[
\sum_j {2^{W_j/C}\over(1+W_j/C)^3}=O(F).          \tag{9}
\]

To see this, write \texttt{x\_j=W\_j/C} and \texttt{f(x)=2\^{}x/(1+x)\^{}3}. On the bounded interval \texttt{{[}0,4{]}}, this function is bounded by an absolute constant. On \texttt{{[}4,infinity)} it is increasing, and

\[
0\le f'(t)\le {\ln2\,2^t\over(1+t)^3}.
\]

Using the tail integral and (1),

\[
\begin{aligned}
\sum_{j:x_j\ge4}f(x_j)
&\le F f(4)+\int_4^\infty f'(t)\#\{j:x_j\ge t\}\,dt\\
&\le F f(4)+{F\ln2\over2}
                 \int_4^\infty {t+3\over(t+1)^3}\,dt
=O(F).
\end{aligned}
\]

The omitted loads below four contribute another \texttt{O(F)}. Multiplying (9) by \texttt{C} supplies an \texttt{O(CF)} initial credit budget. Also \texttt{x\ \textless{}=\ C\_0\ 2\^{}x/(1+x)\^{}3} for all \texttt{x\textgreater{}=0} with an absolute constant \texttt{C\_0}: the ratio is continuous, bounded on compact intervals, and tends to zero as \texttt{x} tends to infinity. Thus this credit dominates the total selected-tree weight up to an absolute factor. For matching phases in which \texttt{LF=O\_a(N)}, this budget is near-linear without an additional logarithmic factor. A generalized-gap phase may have \texttt{F} much larger than its target \texttt{K}, so that inequality must not be assumed there. The stronger mass-sensitive potential in \protect\hyperlink{af-385fd70686-mass-sensitive-tail}{Part D} resolves this issue and gives an \texttt{O(W)} initial budget in terms of total actual demand alone.

As with (8), (9) is a statement about an actual static Conflt output. It is not a claim that this credit is preserved by dynamic repairs.

\hypertarget{af-385fd70686-c.5.-consequence-for-the-previous-chain-example}{%
\subsubsection{C.5. Consequence for the previous chain example}\label{af-385fd70686-c.5.-consequence-for-the-previous-chain-example}}

The whole-component repair example in the historical note \texttt{warm\_start\_whole\_component\_obstruction.md} (identified in the checkpoint recovery inventory) is a valid ledger with a balanced selected forest. However, even its weakened version need not be an output of Conflt. In that family, there are \texttt{Theta(L/k)} heavy cores of normalized load approximately \texttt{k}, while the number of free roots is \texttt{Theta(L)}. Bound (1) permits only \texttt{O(L\ k/2\^{}k)} such Conflt output cores.

For \texttt{k=ceil((log\_2\ log\_2\ L)/4)}, the ratio

\[
\frac{2^k}{k^2}.
\]

tends to infinity. Thus, for sufficiently large parameters, this particular forest provably cannot be a Conflt output, although it satisfies the earlier maximum-load and total-moment inequalities. This is why the chain is not an obstruction to a scheme that retains the actual Conflt history.

\hypertarget{af-385fd70686-c.6.-what-remains-open-dynamically}{%
\subsubsection{C.6. What remains open dynamically}\label{af-385fd70686-c.6.-what-remains-open-dynamically}}

The capped derivative also gives a useful warm-start fact: adding total normalized mass \texttt{p} to arbitrary old sink loads increases the initial \texttt{H\_tau} potential by at most \texttt{gamma\ 2\^{}tau\ p}. No exponential factor depending on the old maximum load is needed.

Nevertheless, (8) accounts only for heavy completed sinks. A later repair that reactivates many previously completed light sinks must restore their full \texttt{H\_tau} contributions, which were not paid by the previous heavy-tail certificate. Repeated recruitment can therefore replenish the same potential. A dynamic proof must bound that replenishment through irreversible history credits, new dirty mass, or erasure charges. No such bound is established here.

\hypertarget{af-385fd70686-part-d.-a-tail-and-initial-budget-proportional-to-actual-demand}{%
\subsection{Part D. A tail and initial budget proportional to actual demand}\label{af-385fd70686-part-d.-a-tail-and-initial-budget-proportional-to-actual-demand}}

\hypertarget{af-385fd70686-d.1.-setup-and-conclusions}{%
\subsubsection{D.1. Setup and conclusions}\label{af-385fd70686-d.1.-setup-and-conclusions}}

Start from a splittable flow with all nonsink and initial sink congestions at most \texttt{C}. Trim it first to exact demands if necessary, using \href{https://github.com/icekylinx/algorithm-frontier/blob/main/manuscripts/graph-backends/exact-demand-flow-trimming.md}{exact demand trimming}. Let the total demand be \texttt{W}, write \texttt{p=W/C}, and let \texttt{x\_j=W\_j/C} be the final normalized sink loads, including deactivated sinks. Thus \texttt{sum\_j\ x\_j=p}.

For every real \texttt{tau\textgreater{}=4}, the Conflt output obeys

\[
\#\{j:x_j\ge\tau\}
\le { (\tau+5)p\over 7\,2^{\tau-1}}.             \tag{1}
\]

Consequently

\[
\sum_{j:x_j\ge\tau}x_j
=O\!\left({p(\tau+1)^2\over2^\tau}\right),       \tag{2}
\]

and

\[
\sum_j {2^{x_j}-1\over(1+x_j)^3}=O(p).           \tag{3}
\]

The credit

\[
C\sum_j {2^{W_j/C}-1\over(1+W_j/C)^3}=O(W)       \tag{4}
\]

dominates the total component weight up to an absolute constant. Unlike a bound \texttt{O(CF)}, this is near-linear in represented weight even when the number \texttt{F} of free bins is much larger than the phase's target improvement.

All conclusions also apply to actual component weights dominated by the guide demands, using an increasing upper envelope for the bounded small-load portion of (3). Exact trimming makes the quantities equal and is the cleanest formulation.

\hypertarget{af-385fd70686-d.2.-a-preliminary-active-potential-with-no-per-sink-baseline}{%
\subsubsection{D.2. A preliminary active potential with no per-sink baseline}\label{af-385fd70686-d.2.-a-preliminary-active-potential-with-no-per-sink-baseline}}

Define the increasing convex function

\[
f(x)=\begin{cases}x,&0\le x\le2,\\2^{x-1},&x\ge2.
\end{cases}                                     \tag{5}
\]

Its right derivative at two is \texttt{2\ ln\ 2\textgreater{}1}, so the join is convex. The sum of (5) over active sinks never increases under Conflt.

Only deactivation needs proof. Normalize the common frontier's flows as \texttt{a,b}, let \texttt{s,t} be the donor and recipient loads, and use the source conditions

\[
0<a\le\min\{1,s\},\qquad a+b\le1,\qquad t\le s+2b.
\]

The donor completes at \texttt{s-a}, and the recipient rises to \texttt{t+a}. Convexity makes \texttt{f(t+a)-f(t)} increasing in \texttt{t}, so it suffices to show

\[
f(s+2-a)-f(s+2-2a)\le f(s).                      \tag{6}
\]

If \texttt{s\textgreater{}=2}, both terms are exponential, and the left side divided by \texttt{f(s)} is \texttt{4(2\^{}\{-a\}-2\^{}\{-2a\})\textless{}=1}.

Suppose \texttt{s\textless{}2}. If the recipient interval lies below two, its increase is \texttt{a\textless{}=s}. If its lower endpoint is at least two, then \texttt{a\textless{}=s/2}. The expression \texttt{2\^{}\{-a\}-2\^{}\{-2a\}} is increasing for \texttt{0\textless{}=a\textless{}=1}, so the increase is at most

\[
2(2^{s/2}-1)\le s.
\]

The last inequality is the chord bound \texttt{2\^{}z\textless{}=1+z} for \texttt{z} in \texttt{{[}0,1{]}}. In the remaining crossing case put \texttt{z=s-a}, where \texttt{0\textless{}z\textless{}1}. Equation (6) becomes exactly \texttt{2\^{}z\textless{}=1+z} after rearrangement.

The initial active sum is \texttt{p}, because all initial loads are at most one. Thus every final load satisfies \texttt{f(x\_j)\textless{}=p}. In particular,

\[
\max_jx_j\le\begin{cases}p,&p\le2,\\1+\log_2p,&p\ge2.\end{cases}\tag{7}
\]

This maximum-load refinement alone does not control all completed loads simultaneously; the next potential does.

\hypertarget{af-385fd70686-d.3.-the-mass-sensitive-capped-slope-potential}{%
\subsubsection{D.3. The mass-sensitive capped-slope potential}\label{af-385fd70686-d.3.-the-mass-sensitive-capped-slope-potential}}

Fix \texttt{tau\textgreater{}=4}, and write

\[
q=\tau+5,\quad D=\tau+7,\quad E=2^{\tau-1},\quad\gamma=7\ln2-1.
\]

Define

\[
H(x)=\begin{cases}
qx,&0\le x\le2,\\
(D-x)2^{x-1},&2\le x\le\tau,\\
E[7+\gamma(x-\tau)],&x\ge\tau.
\end{cases}                                     \tag{8}
\]

This is continuous, increasing, and convex. At two the right derivative is \texttt{2(q\ ln2-1)\textgreater{}=q}, because \texttt{q\textgreater{}=9}. On \texttt{{[}2,tau{]}}, the second derivative is positive since \texttt{D-x\textgreater{}=7}. At \texttt{tau}, the value and derivative match its tangent-line continuation. Its derivative is everywhere at most \texttt{gamma\ E}.

Let \texttt{Psi} be the sum of (8) over active sinks. Aggregation and cycle cancellation preserve sink loads. A nondeactivating balancing move decreases \texttt{Psi} by convexity. We next show the same for deactivation.

\hypertarget{af-385fd70686-d.4.-every-deactivation-is-paid}{%
\subsubsection{D.4. Every deactivation is paid}\label{af-385fd70686-d.4.-every-deactivation-is-paid}}

Use the notation of Section D.2. The active-potential drop is

\[
\Delta=H(s)-[H(t+a)-H(t)].                        \tag{9}
\]

By convexity, replace \texttt{t} in the bracket by its upper bound \texttt{T=s+2-2a}. Write \texttt{T+a=s+2-a}.

\hypertarget{af-385fd70686-d.4.1-donor-load-below-two}{%
\paragraph{D.4.1 Donor load below two}\label{af-385fd70686-d.4.1-donor-load-below-two}}

Suppose \texttt{s\textless{}2}. Since \texttt{a\textless{}=s}, the recipient's largest possible upper endpoint is below four, and hence below \texttt{tau}.

On this entire range write

\[
H(x)=qf(x)-g(x),\qquad
g(x)=\begin{cases}0,&x\le2,\\
                   (x-2)2^{x-1},&2\le x\le\tau.
      \end{cases}                               \tag{10}
\]

The function \texttt{g} is increasing and \texttt{g(s)=0}. The already proved nonincrease of the \texttt{f} potential therefore gives

\[
\Delta\ge q\{f(s)-[f(T+a)-f(T)]\}
                     +g(T+a)-g(T)\ge0.          \tag{11}
\]

This treats the wholly linear, crossing-two, and wholly exponential recipient intervals together.

\hypertarget{af-385fd70686-d.4.2-donor-load-at-least-two}{%
\paragraph{D.4.2 Donor load at least two}\label{af-385fd70686-d.4.2-donor-load-at-least-two}}

If \texttt{s\textgreater{}=tau}, global bounded slope gives

\[
\Delta\ge H(s)-a\gamma E\ge(7-\gamma)E>0.
\]

If \texttt{2\textless{}=s\textless{}tau} and \texttt{T+a\textless{}=tau}, all three values after replacing \texttt{t} by \texttt{T} use the exponential part. The calculation in \protect\hyperlink{af-385fd70686-exponential-moment}{Part B} applies, with exponent scaled by one half and \texttt{D-s\textgreater{}=7\textgreater{}=4}; it gives \texttt{Delta\textgreater{}=2*2\^{}(s-1)\textgreater{}0}.

If \texttt{2\textless{}=s\textless{}tau\textless{}T+a}, then \texttt{s\textgreater{}tau-2+a\textgreater{}=2}. Monotonicity and the bounded derivative give

\[
{\Delta\over E}
>(9-a)2^{a-2}-a\gamma.                           \tag{12}
\]

The right side is decreasing on \texttt{{[}0,1{]}}. Its derivative is increasing before the constant \texttt{gamma} is subtracted, but even its value at one is

\[
\frac{8\ln2-1}{2}-(7\ln2-1)=\frac12-3\ln2<0.
\]

Thus (12) is at least its value at one,

\[
4-\gamma=5-7\ln2>0.
\]

This proves nonincrease in every case.

\hypertarget{af-385fd70686-d.4.3-completed-heavy-sinks-pay-their-entire-receipt}{%
\paragraph{D.4.3 Completed heavy sinks pay their entire receipt}\label{af-385fd70686-d.4.3-completed-heavy-sinks-pay-their-entire-receipt}}

If the donor's completed load \texttt{s-a} is at least \texttt{tau}, both \texttt{s} and \texttt{s-a} lie on the linear tail. Therefore

\[
\Delta\ge H(s)-a\gamma E=H(s-a)\ge7E.            \tag{13}
\]

This charge is stronger than mere nonincrease.

\hypertarget{af-385fd70686-d.5.-summing-the-heavy-load-receipts}{%
\subsubsection{D.5. Summing the heavy-load receipts}\label{af-385fd70686-d.5.-summing-the-heavy-load-receipts}}

Every initial load is at most one, where (8) is linear. Hence

\[
\Psi_{\rm initial}=q\sum_j x_j^{\rm initial}
                   =(\tau+5)p.                  \tag{14}
\]

Use (13) to charge all deactivated sinks whose completed loads are at least \texttt{tau}. The remaining heavy active sinks are still present in the final active potential. Thus

\[
\sum_{j:x_j\ge\tau}H(x_j)\le(\tau+5)p.          \tag{15}
\]

Equations (1) and (2) follow from the linear expression for \texttt{H} on the tail, exactly as in \protect\hyperlink{af-385fd70686-load-tail}{Part C}.

\hypertarget{af-385fd70686-d.6.-a-budget-proportional-to-actual-mass}{%
\subsubsection{D.6. A budget proportional to actual mass}\label{af-385fd70686-d.6.-a-budget-proportional-to-actual-mass}}

Let \texttt{g(x)=(2\^{}x-1)/(1+x)\^{}3}. On \texttt{{[}0,4{]}}, \texttt{g(x)\textless{}=c\ x} for an absolute constant \texttt{c}, so these loads contribute \texttt{O(p)}. On \texttt{{[}4,infinity)}, \texttt{g} is increasing and

\[
0\le g'(t)\le {\ln2\,2^t\over(1+t)^3}.
\]

The tail integral and (1) give

\[
\begin{aligned}
\sum_{j:x_j\ge4}g(x_j)
&\le g(4)\#\{j:x_j\ge4\}
      +\int_4^\infty g'(t)\#\{j:x_j\ge t\}\,dt\\
&\le O(p)+{2p\ln2\over7}
              \int_4^\infty{t+5\over(t+1)^3}\,dt
=O(p).
\end{aligned}
\]

This proves (3). The ratio \texttt{x/g(x)} extends continuously at zero, with limit \texttt{1/ln2}, is bounded on compact intervals, and tends to zero at infinity. Consequently the credit in (4) also dominates the total component mass up to an absolute constant.

\hypertarget{af-385fd70686-d.7.-dynamic-status}{%
\subsubsection{D.7. Dynamic status}\label{af-385fd70686-d.7.-dynamic-status}}

The proof gives an \texttt{O(W)} initial credit budget for an actual static Conflt execution. It avoids an incorrect assumption that \texttt{LF=O(N)} in a generalized-gap phase. However, reactivating previously completed light sinks can still reintroduce potential that was not retained in the heavy receipts (15). The existence of this mass-sensitive budget does not by itself bound how often the same mass can be reprocessed. A dynamic replenishment invariant remains necessary before any lower matching exponent can be claimed.

\hypertarget{af-385fd70686-part-e.-local-rounding-with-fixed-exterior-trees}{%
\subsection{Part E. Local rounding with fixed exterior trees}\label{af-385fd70686-part-e.-local-rounding-with-fixed-exterior-trees}}

\hypertarget{af-385fd70686-e.1.-rejected-strengthening-the-potential-only-covers-active-sinks}{%
\subsubsection{E.1. Rejected strengthening: the potential only covers active sinks}\label{af-385fd70686-e.1.-rejected-strengthening-the-potential-only-covers-active-sinks}}

Write \texttt{W\_T} for the sum of bin weights in a selected tree produced by the deterministic confluent rounding, and let \texttt{F} be the number of free bins at the start of that rounding.

The usual consequence is only the maximum-load statement

\[
W_T\le L(1+\log_2F).
\]

An initially proposed stronger simultaneous inequality was

\[
\sum_T2^{W_T/L}\le2F.\qquad\text{Unsupported proposed strengthening}
\]

This does NOT follow from the published argument. Every initial free sink has demand plus incoming guide flow at most \texttt{L}, and the sum over ACTIVE sinks is at most \texttt{2F}. But a deactivated sink is removed from that sum. Source Lemma 6 separately bounds its final potential by the preceding active potential. This gives an individual maximum-load bound; it does not preserve a sum including all completed sinks.

Therefore an \texttt{O\_a(N)} global exponential recruitment budget cannot be claimed from Conflt without a new theorem. The proposed inference was caught by checking the source's precise definition of its potential and is excluded from all matching bounds.

The earlier chain obstruction in the historical note \texttt{warm\_start\_whole\_component\_obstruction.md} (identified in the checkpoint recovery inventory) remains a valid obstruction to its stated rule. Any stronger restriction on allowable snapshot forests would itself require proof.

\textbf{Subsequent positive result.} The companion \protect\hyperlink{af-385fd70686-exponential-moment}{Part B} supplies such a proved weaker global bound, \texttt{sum\_final\ 2\^{}(W\_T/L)=O(F\ log(F+2))}, using a second potential that explicitly pays for deactivated sinks. It does not reinstate the unsupported \texttt{2F} bound and does not yet establish dynamic replenishment accounting.

\hypertarget{af-385fd70686-e.2.-route-exactly-the-dirty-weights-using-only-dirty-rows}{%
\subsubsection{E.2. Route exactly the dirty weights using only dirty rows}\label{af-385fd70686-e.2.-route-exactly-the-dirty-weights-using-only-dirty-rows}}

Consider any settled bottom ledger and a region \texttt{Q} of bins whose assignments may be repaired. All selected trees outside \texttt{Q} are fixed cores. A bin still carrying an untouched core tag can be treated as outside, even when it also contains some separately credited dirty mass.

For each assigned bin \texttt{v} in \texttt{Q}, use its assigned row's allocation as capacities on arcs from \texttt{v} to the row's neighbors. Delete self-loops and treat every neighbor outside \texttt{Q} as a sink anchor. Free bins in \texttt{Q} are also sinks. The graph has \texttt{O(\textbar{}Q\textbar{})} arcs because each assigned row has at most three neighbors.

The restricted original guide flow has supply

\[
d_v=L-\mathrm{incoming}_{\mathrm{rows\ assigned\ in}\ Q}(v)\ge w_v.
\]

There exists a subflow of this restricted flow that routes EXACTLY \texttt{w\_v} units from every assigned bin \texttt{v} in \texttt{Q} to the sink anchors and free bins, and has no greater arc flow. One proof decomposes the original flow into source-to-sink paths and cycles, discards cycles, and deletes excess paths from each source. An algorithmic construction uses a single exact integral maximum flow with source arcs of capacity \texttt{w\_v}, original arc capacities equal to the restricted allocations, and unrestricted sink arcs. Feasibility follows from the displayed supplies, or directly from the cut inequalities. Its graph has linear size in \texttt{Q}; large integral capacities can use the source framework's capacity-scaling wrapper.

The stronger algorithm in \href{https://github.com/icekylinx/algorithm-frontier/blob/main/manuscripts/graph-backends/exact-demand-flow-trimming.md}{exact demand trimming} removes the need for that maximum-flow call. A directed selected forest performs monotone cycle cancellations and source-to-sink excess removal, with \texttt{O(\textbar{}Q\textbar{})} substantive operations and dynamic-tree overhead. Thus exact dirty demands are obtained in \texttt{O(\textbar{}Q\textbar{}\ polylog\ N)} work.

At a nonsink, own demand plus incoming flow equals outgoing flow and is at most \texttt{L}, since the new flow is coordinatewise bounded by the old row allocation. At a free bin of \texttt{Q}, its own weight plus incoming flow is likewise at most \texttt{L}. Thus this is a supplied flow with nonsink congestion \texttt{C=L} routing the ACTUAL dirty weights, not the potentially larger artificial supplies \texttt{d\_v}.

Identify exterior anchors belonging to the same fixed core as one sink. Give that sink initial local demand equal to the core's current total weight. The original endpoint of each incoming dirty arc is retained as a witness. The sink may now have arbitrarily large initial demand; this does not invalidate Conflt's nonincreasing exponential-potential proof, whose parameter \texttt{C} bounds nonsink congestion. The resulting attachments follow their stored original anchor bins and the unchanged exterior selected tree to its free root.

Consequently a dirty region can be rounded locally, preserving every exterior assignment. If \texttt{Phi\_initial} is the sum of \texttt{2\^{}((old\_core\_load+fractional\_inflow)/L)} before rounding, every output core load is at most \texttt{L\ log\_2\ Phi\_initial}. The active-sink potential never increases; a completed sink is only bounded individually. The work is \texttt{O(\textbar{}Q\textbar{}\ polylog\ N)} with the trimming construction and the near-linear Conflt implementation, at the stated exact integer bit cost.

\hypertarget{af-385fd70686-e.3.-what-remains-open}{%
\subsubsection{E.3. What remains open}\label{af-385fd70686-e.3.-what-remains-open}}

The local lemma does not bound the exponential potential BEFORE rounding: a feasible fractional flow may direct most dirty mass into already heavy exterior cores. The Conflt step cannot undo that initial potential increase without opening additional cores as routing vertices. The old core's internal guide rows may export allocation, so it cannot be replaced by a fictitious capacity-\texttt{L} sink.

A useful continuation would recruit overflowing cores and prove an amortized bound in terms of a genuine global potential, or a bounded replenishment of that potential per gain. Charging a recruited core to its exponential weight is insufficient on its own: later local rounding may recreate that potential at another core, permitting the same mass to be charged repeatedly.

No such decrease or replenishment theorem has been established here. In particular, a three-level cost expression such as \texttt{N\ b1\ +\ N\ b2/b1\ +\ N\ K/b2} is only exponent arithmetic until this local recruitment/repair amortization is proved.

\hypertarget{af-385fd70686-source-and-restoration-record}{%
\subsection{Source and restoration record}\label{af-385fd70686-source-and-restoration-record}}

These results were already present in the 24 October 2026 uniform checkpoint. This page restores their full arguments in a semantic location; it does not claim that they were discovered during the repository reorganization. The original UTF-8 bytes were recovered and checked against the archive inventory. See the \href{https://github.com/icekylinx/algorithm-frontier/blob/main/provenance/checkpoint-restoration.md}{checkpoint restoration record}.

\begin{longtable}[]{@{}
  >{\raggedright\arraybackslash}p{(\columnwidth - 2\tabcolsep) * \real{0.5000}}
  >{\raggedright\arraybackslash}p{(\columnwidth - 2\tabcolsep) * \real{0.5000}}@{}}
\toprule\noalign{}
\begin{minipage}[b]{\linewidth}\raggedright
Original note
\end{minipage} & \begin{minipage}[b]{\linewidth}\raggedright
SHA-256 of recovered source
\end{minipage} \\
\midrule\noalign{}
\endhead
\bottomrule\noalign{}
\endlastfoot
\texttt{confluent\_rounding\_acceleration.md} & \texttt{1a7daa21a9560431438d57a21ad0bc24aedecc3f7dc46a7c53e25089c35d4f86} \\
\texttt{confluent\_exponential\_mass.md} & \texttt{e729d94faed0113e28ba0291804050f1d4c6417b2594697d64fc157e70d4cda3} \\
\texttt{confluent\_load\_tail.md} & \texttt{9e89aca010eb5f02395c60c12fde089ed3ae7b2b1a45ef77dc8a9e30fcb9fe41} \\
\texttt{confluent\_mass\_sensitive\_tail.md} & \texttt{88ff5e7a7b10ebe02f42e266f891ac590090b149b55ff5833c62e0a5eb25b151} \\
\texttt{warm\_start\_potential\_extension.md} & \texttt{53516c069d953f29d03cd6c080aa4a60dcee2aa2648b483399c22dc777c0f72e} \\
\end{longtable}
